We now pose a mathematical problem whose solutions will describe the equilibrium configurations of the body. A material whose equilibria are described by such a problem is said to be \emph{elastic}. We begin by defining an \emph{energy functional}:
\begin{equation}\label{eq:elasticity-total-energy-functional}
\mathcal{E}(\bm{p}):=\int_\Omega\mathcal{W}(\bm{x},\bm{F}(\bm{x}))\,\mathrm{d}\bm{x}+\int_\Omega\Psi(\bm{p}(\bm{x}))\,\mathrm{d}\bm{x}.
\end{equation}
Here $\bm{F}:=\nabla\bm{p}$ and $\mathcal{W}:\Omega\times Q\to\mathbb{R}$ is the \emph{stored energy density}, where $Q$ is the set of $3\times3$ matrices with positive determinant. The function $\Psi:\mathbb{R}^3\to\mathbb{R}$ is the \emph{potential energy density}. The first term in the energy is called the \emph{stored energy functional}, and this describes the energy stored from mechanical deformation within the material. The second term is a \emph{potential energy functional} and it is the energy from exterior forces (assumed to be conservative).

We make the following assumptions about the density functionals.\begin{enumerate}\item We assume that $\mathcal{W}\in C^2(\Omega\times Q)$ and $\Psi\in C^1(\mathbb{R}^3)$.\item For some $p>3$, there exists a constant $k>0$ and a function $\omega\in C_b(\Omega)$ such that
\begin{equation}\label{eq:elasticity-stored-energy-p-growth-bound}
\mathcal{W}(\bm{x},\bm{F})\geq\omega(\bm{x})+k|\bm{F}|^p.
\end{equation}
Note: From now on this $p$ is to be used in the definition of $\mathcal{D}_{\mathcal{E}}$.\item There exists a constant $C$ such that for every $\bm{p}\in\mathcal{D}_{\mathcal{E}}$ we have
\begin{equation}\label{eq:elasticity-potential-energy-uniform-lower-bound}
\int_\Omega\Psi(\bm{p}(\bm{x}))\,\mathrm{d}\bm{x}\geq C.
\end{equation}
\item There exists $\bm{p}\in\mathcal{D}_{\mathcal{E}}$ such that $\mathcal{E}(\bm{p})<\infty$.\item For every $\bm{x}\in\Omega$, we have $\mathcal{W}(\bm{x},\bm{F})\to\infty$ as $\det F\to0$.\end{enumerate}

We now define the equilibrium configurations of the body to be those that minimize the energy $\mathcal{E}$; i.e., we say that $\overline{\bm{p}}\in\mathcal{D}_{\mathcal{E}}$ is an \emph{equilibrium state} or a \emph{minimizer} if
\begin{equation}\label{eq:elasticity-equilibrium-global-minimization}
\mathcal{E}(\overline{\bm{p}})\leq\mathcal{E}(\bm{p})\quad\text{for all }\bm{p}\in\mathcal{D}_{\mathcal{E}}.
\end{equation}
The primary goal of this section is to examine the question of existence of solutions of the minimization problem \eqref{eq:elasticity-equilibrium-global-minimization}. But before doing so, we wish to expose the connection of this problem to elliptic partial differential equations. We do so in the following result.
